\part{KIẾN THỨC CỐT LÕI VÀ THỰC HÀNH}

\chapter{Thiết lập môi trường và Cấu trúc dự án}

\section{Công cụ và Môi trường phát triển}
Để xây dựng và triển khai dự án Trello Mini đáp ứng được các tiêu chuẩn khắt khe của ngành công nghiệp phần mềm hiện đại, nhóm đã thiết lập một hệ sinh thái công cụ phát triển tối ưu nhất:
\begin{itemize}
    \item \textbf{Môi trường Runtime (Runtime Environment):} Node.js phiên bản LTS (Long Term Support) v20+. Node.js không chỉ là môi trường thực thi mà còn cung cấp \texttt{npm} (Node Package Manager), công cụ cốt lõi giúp quản lý, cài đặt và theo dõi các thư viện phụ thuộc (dependencies) trong suốt vòng đời của dự án.
    \item \textbf{Môi trường lập trình (IDE):} Visual Studio Code (VSCode). Đây là sự lựa chọn số một của các lập trình viên Web nhờ dung lượng nhẹ, hiệu năng cao và hệ sinh thái Extensions khổng lồ. Nhóm đã cấu hình sẵn một bộ quy chuẩn cho VSCode bao gồm: \textit{ESLint} (kiểm tra các lỗi cú pháp và logic mã nguồn), \textit{Prettier} (đảm bảo format code đồng nhất cho mọi thành viên), và \textit{Tailwind CSS IntelliSense} (hỗ trợ tự động gợi ý class CSS, tăng tốc độ cắt HTML/CSS lên 300\%).
\end{itemize}

\section{Kiến trúc Siêu tốc của Vite: Deep Dive}
Thay vì sử dụng các công cụ đóng gói truyền thống như Webpack vốn đã trở nên già cỗi và chậm chạp, dự án lựa chọn \textbf{Vite}. Để hiểu tại sao Vite lại là "kẻ thay đổi cuộc chơi", ta cần đi sâu vào kiến trúc bên trong của nó.

\subsection{ESBuild vs Rollup: Sự kết hợp hoàn hảo}
Vite sử dụng hai trình đóng gói khác nhau cho hai mục đích khác nhau, tối ưu hóa cả quá trình phát triển lẫn sản phẩm đầu ra:
\begin{itemize}
    \item \textbf{Development (ESBuild):} Trong quá trình phát triển, Vite sử dụng \texttt{esbuild} để pre-bundle các thư viện phụ thuộc (dependencies). \texttt{esbuild} được viết bằng Go và có tốc độ nhanh hơn từ 10 đến 100 lần so với các công cụ dựa trên JavaScript. Nó tận dụng tối đa đa luồng (multi-threading) của CPU, giúp server khởi động gần như tức thì ngay cả với hàng trăm dependency.
    \item \textbf{Production (Rollup):} Đối với bản build cuối cùng, Vite tin dùng \texttt{Rollup}. Mặc dù chậm hơn ESBuild, Rollup cung cấp khả năng tối ưu hóa mã nguồn cực cao như:
    \begin{itemize}
        \item \textit{Tree-shaking:} Loại bỏ các đoạn code không bao giờ được dùng đến (dead code).
        \item \textit{Scope Hoisting:} Gộp các module lại để giảm kích thước file và tăng tốc độ thực thi.
        \item \textit{Advanced Chunking:} Chia nhỏ file thông minh để tận dụng cơ chế cache của trình duyệt.
    \end{itemize}
\end{itemize}

\subsection{Native ESM: Giải mã sự khác biệt với Webpack}
Webpack hoạt động theo cơ chế "Bundle-based": Nó phải quét toàn bộ ứng dụng, xây dựng cây phụ thuộc (Dependency Graph) và đóng gói tất cả thành một file khổng lồ (bundle) trước khi có thể hiển thị trang web. Khi dự án lớn lên, thời gian đợi chờ này tăng theo cấp số nhân ($O(n)$).

Vite giải quyết bài toán này bằng cách sử dụng \textbf{Native ESM (ES Modules)} của trình duyệt. Thay vì đóng gói, Vite chỉ đơn giản là phục vụ các file nguồn theo yêu cầu của trình duyệt. 
\begin{enumerate}
    \item Trình duyệt gặp lệnh \texttt{import}.
    \item Trình duyệt gửi một request HTTP đến server của Vite.
    \item Vite nhận request, biên dịch đúng file đó sang JavaScript tương thích (nếu cần) và gửi lại.
\end{enumerate}
Nhờ đó, tốc độ khởi động của Vite là hằng số ($O(1)$), không phụ thuộc vào độ lớn của dự án.

\subsection{Hot Module Replacement (HMR) và Hiệu quả lập trình}
Vite thực hiện HMR qua Native ESM một cách cực kỳ tinh tế. Khi một file thay đổi, Vite chỉ vô hiệu hóa module đó và các module liên quan trực tiếp. Trình duyệt chỉ cần tải lại đúng file đó mà không mất đi trạng thái (state) hiện tại của ứng dụng. Điều này giúp lập trình viên thấy ngay kết quả thay đổi sau khi nhấn Ctrl+S chỉ trong vài mili giây.

\section{Kiến trúc mã nguồn Feature-Sliced Design (FSD)}
Dự án áp dụng \textbf{Feature-Sliced Design (FSD)} để giải quyết bài toán bảo trì lâu dài. FSD không chỉ là cách đặt tên thư mục, nó là một tư duy phân cấp logic:

\begin{table}[h]
\centering
\begin{tabular}{|l|l|l|}
\hline
\textbf{Layer} & \textbf{Trách nhiệm} & \textbf{Ví dụ trong dự án} \\ \hline
App & Cấu hình toàn cục & \texttt{router.tsx}, \texttt{main.tsx} \\ \hline
Pages & Giao diện các trang & \texttt{DashboardPage}, \texttt{BoardPage} \\ \hline
Features & Logic nghiệp vụ & \texttt{AddCardForm}, \texttt{DeleteBoardButton} \\ \hline
Entities & Logic dữ liệu & \texttt{useBoardStore}, \texttt{ICard interface} \\ \hline
Shared & Thành phần dùng chung & \texttt{Button.tsx}, \texttt{api.ts} \\ \hline
\end{tabular}
\caption{Phân tầng kiến trúc FSD trong Trello Mini}
\end{table}

\chapter{Hệ thống Type tiên tiến với TypeScript}

\section{Generics: Tái sử dụng mã nguồn an toàn}
Generics là tính năng quan trọng nhất để tạo ra các component và hàm linh hoạt. Trong ứng dụng Kanban, chúng ta thường xuyên phải xử lý các danh sách (Bảng, Cột, Thẻ).

\begin{Verbatim}
// Interface cho một API Response chung
interface ApiResponse<T> {
  data: T;
  status: number;
  message: string;
}

// Cách dùng:
const boardResponse: ApiResponse<IBoard> = { ... };
const cardsResponse: ApiResponse<ICard[]> = { ... };
\end{Verbatim}
Việc sử dụng \texttt{<T>} giúp chúng ta định nghĩa cấu trúc dữ liệu một lần và sử dụng cho nhiều loại thực thể khác nhau mà vẫn giữ được sự kiểm soát chặt chẽ của TypeScript.

\section{Utility Types: Biến đổi linh hoạt}
\begin{itemize}
    \item \textbf{Partial<T>}: Cực kỳ hữu ích trong các form cập nhật. Khi người dùng chỉ muốn sửa tên thẻ, chúng ta truyền một \texttt{Partial<ICard>} vào store thay vì bắt buộc phải truyền toàn bộ đối tượng.
    \item \textbf{Required<T>}: Ngược lại với Partial, bắt buộc tất cả thuộc tính phải có giá trị.
    \item \textbf{Readonly<T>}: Ngăn chặn việc thay đổi dữ liệu trực tiếp, ép buộc lập trình viên phải sử dụng các hàm cập nhật state của Zustand (tuân thủ nguyên tắc Immutability).
\end{itemize}

\section{Discriminated Unions và Pattern Matching}
Trong logic kéo thả, chúng ta cần phân biệt đâu là Thẻ (Card) và đâu là Cột (List) để có cách xử lý va chạm khác nhau.

\begin{Verbatim}
type DragItem = 
  | { type: 'CARD'; data: ICard; listId: string }
  | { type: 'LIST'; data: IList; boardId: string };

function onDragStart(item: DragItem) {
  if (item.type === 'CARD') {
    // TypeScript biết chắc chắn ở đây item có listId
    console.log("Đang kéo thẻ từ cột:", item.listId);
  } else {
    // Ở đây item chỉ có boardId
    console.log("Đang kéo cột trong bảng:", item.boardId);
  }
}
\end{Verbatim}

\chapter{React 18/19 và Hiệu năng Đột phá}

\section{Concurrent Mode: Xử lý đa nhiệm trong UI}
React 18 không chỉ là một bản cập nhật thông thường, nó là sự thay đổi về mô hình thực thi. Với \textbf{Concurrent Rendering}, React có thể chuẩn bị giao diện mới ở "hậu trường" mà không làm gián đoạn các tương tác hiện tại của người dùng. 

\section{Deep Dive into useTransition}
Khi người dùng gõ từ khóa tìm kiếm thẻ, hai việc xảy ra:
\begin{enumerate}
    \item Ô input cần hiển thị ký tự người dùng vừa gõ (Ưu tiên 1 - Phải ngay lập tức).
    \item Danh sách thẻ bên dưới cần được lọc lại (Ưu tiên 2 - Có thể trễ 100-200ms).
\end{enumerate}
Nếu không có \texttt{useTransition}, React sẽ cố gắng làm cả hai cùng lúc. Nếu danh sách thẻ quá lớn, ô input sẽ bị khựng (lag). \texttt{useTransition} cho phép React "trì hoãn" việc lọc danh sách để ưu tiên hiển thị ô input.

\section{Suspense và Error Boundaries}
Dự án kết hợp \textbf{Suspense} để hiển thị Skeleton Screen (giao diện khung xương) trong khi đợi dữ liệu từ LocalStorage hoặc API, mang lại cảm giác ứng dụng phản hồi nhanh hơn thực tế (Perceived Performance). Đi kèm với đó là \textbf{Error Boundaries} để đảm bảo nếu một thẻ bị lỗi render, toàn bộ bảng Kanban vẫn hoạt động bình thường thay vì bị trắng trang.

\chapter{Zustand: Quản lý State Hiện đại}

\section{Tại sao không dùng Redux hay Context API?}
\begin{itemize}
    \item \textbf{Context API:} Quá đơn giản cho các ứng dụng phức tạp. Mỗi khi một giá trị nhỏ trong Context thay đổi, toàn bộ các component con đều bị render lại, dẫn đến hiệu năng cực kém khi bảng Kanban có hàng trăm thẻ.
    \item \textbf{Redux Toolkit:} Quá rườm rà. Để thêm một tính năng "Xóa thẻ", ta cần viết Action, Reducer, Slice, Dispatch. Mã nguồn trở nên nặng nề không cần thiết.
    \item \textbf{Zustand:} Cung cấp cơ chế \textit{Selectors} giống Redux giúp tối ưu render, nhưng lại có cú pháp khai báo cực kỳ ngắn gọn giống như các Hook bình thường.
\end{itemize}

\section{Zustand Middleware và Data Persistence}
\begin{Verbatim}
export const useBoardStore = create<StoreState>()(
  devtools(
    persist(
      immer((set) => ({ ... })),
      { name: 'trello-storage' }
    )
  )
);
\end{Verbatim}
Việc lồng ghép các Middleware cho thấy sức mạnh của Zustand:
\begin{itemize}
    \item \textbf{Immer:} Biến code cập nhật mảng lồng nhau từ 10 dòng thành 1 dòng.
    \item \textbf{Persist:} Tự động đồng bộ với LocalStorage, giải quyết bài toán lưu trữ dữ liệu phía client.
    \item \textbf{Devtools:} Cho phép debug trạng thái ứng dụng một cách chuyên nghiệp.
\end{itemize}

\chapter{Tailwind CSS: Hệ thống Giao diện Linh hoạt}

\section{Sức mạnh của JIT (Just-In-Time)}
Trong dự án này, Tailwind không chỉ là một thư viện CSS, nó là một \textit{Engine}. Nó quét toàn bộ mã nguồn React, tìm các class bạn đã dùng và chỉ sinh ra đúng bấy nhiêu CSS. Kết quả là file CSS gửi đến người dùng luôn ở mức tối thiểu, giúp tốc độ tải trang đạt điểm tuyệt đối trên Google PageSpeed Insights.

\section{Atomic Design và Reusability}
Chúng ta không viết CSS cho từng trang. Thay vào đó, chúng ta xây dựng bộ thư viện UI:
\begin{itemize}
    \item \texttt{Button.tsx}: Chứa mọi logic về màu sắc, kích thước, trạng thái loading.
    \item \texttt{Input.tsx}: Chứa logic về focus, border, error states.
\end{itemize}
Khi cần thay đổi màu sắc chủ đạo của toàn bộ ứng dụng, chúng ta chỉ cần sửa đúng một dòng trong file \texttt{index.css} tại khối \texttt{@theme}.

\chapter{Tối ưu hóa Hiệu năng và Trải nghiệm người dùng}

\section{Kỹ thuật Virtualization (Tầm nhìn tương lai)}
Mặc dù dự án hiện tại xử lý tốt hàng trăm thẻ, nhưng với hàng ngàn thẻ, chúng ta cần đến \textbf{Virtualization} (chỉ render những thẻ đang hiển thị trên màn hình). Đây là hướng phát triển tiếp theo của dự án để đảm bảo hiệu năng không giới hạn.

\section{Debouncing và Throttling trong Xử lý sự kiện}
Các hành động như thay đổi kích thước cửa sổ hoặc cuộn trang được tối ưu bằng \textbf{Throttling} để giảm số lần tính toán lại vị trí các thẻ, giúp CPU không bị quá tải.

\section{Accessibility (A11y)}
Dự án không chỉ chú trọng vẻ đẹp mà còn cả tính tiếp cận. Các thẻ và nút đều có thuộc tính \texttt{aria-label}, hỗ trợ người khiếm thị sử dụng Screen Readers để tương tác với bảng Kanban một cách dễ dàng.
